% !TEX root = ../main.tex
%% Corrigés — Chapitre 10 — Fonction exponentielle
%% Tous les résultats ont été vérifiés par calcul formel (sympy).

\corrigesChapitre{10}{Fonction exponentielle}

\begin{deuxcolonnes}

\begin{corrige}{1}{Propriétés algébriques et calculs de base}
\begin{questions}
  \item \methode{$e^{a}\times e^{b}=e^{a+b}$}\\
    $e^{2x+3}\times e^{-x}=e^{(2x+3)+(-x)}=\res{e^{x+3}}$
  \item \methode{$e^{a}/e^{b}=e^{a-b}$}\\
    $\dfrac{e^{3x}}{e^{2x+1}}=e^{3x-(2x+1)}=\res{e^{x-1}}$
  \item \methode{$(e^{a})^{n}=e^{na}$}\\
    $(e^{x})^{2}\times e^{-3x}=e^{2x}\times e^{-3x}=\res{e^{-x}}$\\
    soit $a=-1$ et $b=0$.
  \item $e^{x}\times e^{-x}=e^{x+(-x)}=e^{0}=\res{1}$
  \item $e^{x+2}\times e^{3x-1}=e^{4x+1}$, donc\\
    $\dfrac{e^{4x+1}}{e^{2x}}=e^{(4x+1)-2x}=\res{e^{2x+1}}$
\end{questions}
\end{corrige}

\begin{corrige}{2}{Limites et variations de base}
\begin{questions}
  \item \methode{limite de référence}\\
    $\displaystyle\lim_{x\to+\infty}e^{x}=\res{+\infty}$
  \item \methode{limite de référence}\\
    $\displaystyle\lim_{x\to-\infty}e^{x}=\res{0}$
  \item \methode{taux d'accroissement en $0$}\\
    $\dfrac{e^{x}-1}{x}=\dfrac{e^{x}-e^{0}}{x-0}$ est le taux d'accroissement de $\exp$ en $0$, donc\\
    $\displaystyle\lim_{x\to0}\frac{e^{x}-1}{x}=\exp'(0)=e^{0}=\res{1}$
  \item La fonction exponentielle est strictement positive :\\
    pour tout réel $x$, $\res{e^{x}>0}$
  \item $2<3$ et $\exp$ est strictement croissante sur $\R$, donc $\res{e^{2}<e^{3}}$
\end{questions}
\end{corrige}

\begin{corrige}{3}{Dérivées de base}
\methode{si $u$ est dérivable, $(e^{u})'=u'e^{u}$}
\begin{questions}
  \item $\exp$ est égale à sa dérivée :\\ $f'(x)=\res{e^{x}}$
  \item $u(x)=3x$ et $u'(x)=3$, donc\\ $g'(x)=\res{3e^{3x}}$
  \item $u(x)=-2x+1$ et $u'(x)=-2$, donc\\ $h'(x)=\res{-2e^{-2x+1}}$
  \item $k'(x)=5\times 2e^{2x}=\res{10e^{2x}}$
  \item $u(x)=x^{2}$ et $u'(x)=2x$, donc\\ $m'(x)=\res{2xe^{x^{2}}}$
\end{questions}
\end{corrige}

\begin{corrige}{4}{Équations et inéquations}
\begin{questions}
  \item \methode{$e^{a}=e^{b}\iff a=b$}\\
    $2x+1=x-3\iff x=-4$\\
    $\res{S=\ens{-4}}$
  \item \methode{pour $k>0$ :}\\
    $e^{x}=k\iff x=\ln k$\\
    $5>0$, donc $\res{S=\ens{\ln 5}}$\enspace($\ln5\approx1{,}609$)
  \item \methode{changement de variable}\\
    On pose $X=e^{x}$, avec $X>0$ :\\
    $X^{2}-4X+3=0$ et $\Delta=4$,\\
    d'où $X=1$ ou $X=3$.\\
    $e^{x}=1\iff x=0$\enspace et\\ $e^{x}=3\iff x=\ln3$\\
    $\res{S=\ens{0\pv\ln3}}$
  \item \methode{$\exp$ est strictement croissante}\\
    $x-1<2x+3\iff x>-4$\\
    $\res{S=\intOO{-4}{+\infty}}$
  \item $7>0$ et $\exp$ est strictement croissante :\\
    $e^{x}\geq7\iff x\geq\ln7$\\
    $\res{S=\intFO{\ln7}{+\infty}}$\enspace($\ln7\approx1{,}946$)
\end{questions}
\end{corrige}

\begin{corrige}{5}{Équations et inéquations avancées}
\methode{on pose $X=e^{x}>0$}
\begin{questions}
  \item $X^{2}-5X+6=(X-2)(X-3)$, donc $e^{x}=2$ ou $e^{x}=3$ :\\ $\res{S=\ens{\ln2\pv\ln3}}$
  \item $x^{2}=3x-2\iff(x-1)(x-2)=0$ :\\ $\res{S=\ens{1\pv2}}$
  \item $(X-1)(X-2)\leq0\iff1\leq X\leq2$,\\ soit $0\leq x\leq\ln2$ : $\res{S=\intFF{0}{\ln2}}$
  \item On divise par $e^{x-1}>0$ : l'inéquation équivaut à $e^{2}>2$, toujours vraie : $\res{S=\R}$
  \item $e^{x^{2}-4x+3}\geq e^{0}\iff(x-1)(x-3)\geq0$ :\\ $\res{S=\intOF{-\infty}{1}\cup\intFO{3}{+\infty}}$
\end{questions}
\end{corrige}

\begin{corrige}{6}{Exponentielle et logarithme}
\begin{questions}
  \item Il faut $e^{2x}>1$, soit $x>0$. Puis $e^{2x}=e^{2}+1$ :\\ $\res{S=\ens{\frac12\ln\left(e^{2}+1\right)}}$
  \item Il faut $x>-1$ ; alors $e^{\ln(x+1)}=x+1=3$ :\\ $\res{S=\ens{2}}$
  \item Il faut $x^{2}-1>0$ et $x+1>0$, soit $x>1$. Puis $x^{2}-1=x+1\iff(x+1)(x-2)=0$ ; seul $2$ convient : $\res{S=\ens{2}}$
  \item Il faut $x>0$ ; alors $e^{2\ln x}=x^{2}$ et $x^{2}-3x+2=0$ :\\ $\res{S=\ens{1\pv2}}$
  \item Il faut $e^{x}>1$, soit $x>0$. Puis $e^{x}-1\geq e$, soit $x\geq\ln(e+1)$ :\\ $\res{S=\intFO{\ln(e+1)}{+\infty}}$
\end{questions}
\end{corrige}

\end{deuxcolonnes}

%% ---------------------------------------------------------------------
\begin{corrige}{7}{Limites et croissances comparées}
\begin{questions}
  \item $f'(x)=\res{e^{x}-1}$.\enspace
    Pour $g$, on utilise $\left(\dfrac{u}{v}\right)'=\dfrac{u'v-uv'}{v^{2}}$ avec $u(x)=e^{x}$ et $v(x)=x$ :
    \[ g'(x)=\frac{e^{x}\times x-e^{x}\times1}{x^{2}}=\res{\frac{e^{x}(x-1)}{x^{2}}} \]
  \item Pour $f$ : $f'(x)=0\iff e^{x}=1\iff x=0$, avec $f'(x)<0$ si $x<0$ et $f'(x)>0$ si $x>0$.\\
    Donc $f$ est décroissante sur $\intOF{-\infty}{0}$ et croissante sur $\intFO{0}{+\infty}$.\\[1mm]
    Pour $g$ : sur $\intOO{0}{+\infty}$, $e^{x}>0$ et $x^{2}>0$, donc $g'(x)$ est du signe de $x-1$.\\
    Donc $g$ est décroissante sur $\intOF{0}{1}$ et croissante sur $\intFO{1}{+\infty}$.
  \item \methode{lever la forme indéterminée « $\infty-\infty$ »}\\
    En $+\infty$, on factorise : $f(x)=e^{x}\left(1-\dfrac{x}{e^{x}}\right)$. Par croissance comparée :
    \[ \lim_{x\to+\infty}\frac{x}{e^{x}}=0 \qquad\text{donc}\qquad \lim_{x\to+\infty}f(x)=\res{+\infty} \]
    En $-\infty$ :
    \[ \lim_{x\to-\infty}e^{x}=0 \ \text{ et } \lim_{x\to-\infty}(-x)=+\infty
       \qquad\text{donc}\qquad \lim_{x\to-\infty}f(x)=\res{+\infty} \]
  \item En $0^{+}$ :
    \[ \lim_{x\to0^{+}}e^{x}=1 \ \text{ et } \lim_{x\to0^{+}}x=0^{+}
       \qquad\text{donc}\qquad \lim_{x\to0^{+}}g(x)=\res{+\infty} \]
    En $+\infty$, par croissance comparée :
    \[ \lim_{x\to+\infty}\frac{e^{x}}{x}=\res{+\infty} \]
  \item On a $f(0)=e^{0}-0=1$ et $g(1)=\dfrac{e^{1}}{1}=e$.\par\vspace{2mm}
    \noindent
    \begin{minipage}[t]{0.49\linewidth}\centering
      \begin{tableau}
        \otTab[1.8]{$x$/0.9,$f'(x)$/0.9,$f$/1.9}{$-\infty$,$0$,$+\infty$}
        \tkzTabLine{,-,z,+,}
        \tkzTabVar{+/$+\infty$,-/$1$,+/$+\infty$}
      \end{tableau}
    \end{minipage}\hfill
    \begin{minipage}[t]{0.49\linewidth}\centering
      \begin{tableau}
        \otTab[1.8]{$x$/0.9,$g'(x)$/0.9,$g$/1.9}{$0$,$1$,$+\infty$}
        \tkzTabLine{d,-,z,+,}
        \tkzTabVar{D+/ /$+\infty$,-/$e$,+/$+\infty$}
      \end{tableau}
    \end{minipage}
\end{questions}
\end{corrige}

%% ---------------------------------------------------------------------
\begin{corrige}{8}{Dérivée et applications}
\begin{questions}
  \item \methode{$(uv)'=u'v+uv'$, avec $u(x)=x$ et $v(x)=e^{x}$}\\
    $g'(x)=1\times e^{x}+x\times e^{x}=\res{e^{x}(1+x)}$
  \item Comme $e^{x}>0$ pour tout réel $x$, $g'(x)$ est du signe de $1+x$ :
    $g'(x)<0$ si $x<-1$, $g'(-1)=0$ et $g'(x)>0$ si $x>-1$.
  \item $g(-1)=(-1)\times e^{-1}=-\dfrac{1}{e}$.
    \begin{center}
      \begin{tableau}
        \otTab[3]{$x$/0.9,$g'(x)$/0.9,$g$/1.9}{$-\infty$,$-1$,$+\infty$}
        \tkzTabLine{,-,z,+,}
        \tkzTabVar{+/$0$,-/$-\frac{1}{e}$,+/$+\infty$}
      \end{tableau}\\[1mm]
      \legende{Les limites sont justifiées à la question 4.}
    \end{center}
    $g$ est décroissante sur $\intOF{-\infty}{-1}$, croissante sur $\intFO{-1}{+\infty}$,
    et admet un minimum égal à $-\frac{1}{e}$ en $x=-1$.
\end{questions}

\noindent
\begin{minipage}[t]{0.54\linewidth}\vspace{0pt}
\begin{questions}[start=4]
  \item En $+\infty$ : $x$ et $e^{x}$ tendent vers $+\infty$, donc, par produit :
    \[ \lim_{x\to+\infty}g(x)=\res{+\infty} \]
    En $-\infty$, par croissance comparée :
    \[ \lim_{x\to-\infty}xe^{x}=\res{0} \]
  \item \methode{$y=g'(a)(x-a)+g(a)$, avec $a=0$}\\
    $g(0)=0$ et $g'(0)=e^{0}(1+0)=1$.\\
    La tangente $T$ en $0$ a pour équation $\res{y=x}$.
\end{questions}
\end{minipage}\hfill
\begin{minipage}[t]{0.42\linewidth}\vspace{0pt}
\centering
\begin{tikzpicture}
\begin{axis}[otGraphe, width=\linewidth,
  xmin=-4.3, xmax=1.4, ymin=-0.9, ymax=2.4,
  xtick={-4,-3,-2,-1}, ytick={1,2},
  xticklabel style={anchor=south,yshift=1.5pt}]
  \addplot[otGreen,line width=0.9pt,domain=-0.9:1.4] {x};
  \addplot[otMuted,dashed,line width=0.5pt] coordinates {(-1,0) (-1,-0.3679)};
  \addplot[otNavy,line width=1.2pt,domain=-4.3:0.93,samples=200,smooth] {x*exp(x)};
  \addplot[only marks,mark=*,mark size=1.5pt,otNavy] coordinates {(-1,-0.3679) (0,0)};
  \node[otEtiquette,anchor=north] at (axis cs:-1,-0.42) {$\left(-1\pv-\frac{1}{e}\right)$};
  \node[anchor=west,font=\small,text=otNavy] at (axis cs:0.97,1.95) {$\mathcal{C}_g$};
  \node[anchor=west,font=\small,text=otGreen!70!otInk] at (axis cs:1.0,0.75) {$T$};
\end{axis}
\end{tikzpicture}\\[1mm]
\legende{Courbe de $g$ et sa tangente $T$ au point d'abscisse $0$.}
\end{minipage}
\end{corrige}

%% ---------------------------------------------------------------------
\begin{corrige}{9}{Modélisation}
\begin{questions}
  \item $\res{N(t)=1\,000\,e^{0{,}02t}}$ et $N(50)=1\,000\,e\approx\res{2\,718\text{ bactéries}}$.
  \item $e^{0{,}02t}=2\iff0{,}02\,t=\ln2\iff t=50\ln2\approx\res{34{,}7\text{ h}}$.
  \item $N'(t)=1\,000\times0{,}02\,e^{0{,}02t}=20\,e^{0{,}02t}=0{,}02\,N(t)$ : la vitesse de croissance est, à chaque instant,
    proportionnelle à l'effectif (2\,\% de l'effectif par heure).
  \item $e^{0{,}02t}=10\iff t=50\ln10\approx\res{115{,}1\text{ h}}$.
\end{questions}
\end{corrige}


%% ---------------------------------------------------------------------
\begin{corrige}{10}{Étude complète et applications}
\begin{questions}
  \item \methode{$(e^{u})'=u'e^{u}$, avec $u(x)=2x$}\\
    $f$ est dérivable sur $\R$, avec $(e^{2x})'=2e^{2x}$ et $(-3e^{x})'=-3e^{x}$ ; la dérivée d'une constante est nulle. Donc, pour tout réel $x$ :
    \[ f'(x)=\res{2e^{2x}-3e^{x}} \]
  \item On factorise : $f'(x)=e^{x}\left(2e^{x}-3\right)$. Comme $e^{x}>0$ pour tout réel $x$, $f'(x)$ est du signe de $2e^{x}-3$ :
    \[ 2e^{x}-3>0\iff e^{x}>\frac{3}{2}\iff x>\ln\frac{3}{2}. \]
    \begin{center}
      \begin{tableau}
        \otTab[3.2]{$x$/0.9,$f'(x)$/0.9,$f$/1.9}{$-\infty$,$\ln\frac{3}{2}$,$+\infty$}
        \tkzTabLine{,-,z,+,}
        \tkzTabVar{+/$2$,-/$-\frac{1}{4}$,+/$+\infty$}
      \end{tableau}\\[1mm]
      \legende{Les valeurs du tableau sont justifiées aux questions 3 et 6.}
    \end{center}
    $f$ est décroissante sur $\intOF{-\infty}{\ln\frac{3}{2}}$ et croissante sur $\intFO{\ln\frac{3}{2}}{+\infty}$.
  \item D'après la question 2, $f$ admet un minimum en $x=\ln\frac{3}{2}$. Comme $e^{\ln\frac{3}{2}}=\frac{3}{2}$ :
    \[ f\Big(\ln\frac{3}{2}\Big)=\Big(\frac{3}{2}\Big)^{2}-3\times\frac{3}{2}+2
       =\frac{9}{4}-\frac{18}{4}+\frac{8}{4}=\res{-\frac{1}{4}} \]
\end{questions}

\noindent
\begin{minipage}[t]{0.54\linewidth}\vspace{0pt}
\begin{questions}[start=4]
  \item \methode{changement de variable}\\
    On pose $X=e^{x}$, avec $X>0$ : $X^{2}-3X+2=0$.\\
    $\Delta=9-8=1$, d'où $X=1$ ou $X=2$.\\
    $e^{x}=1\iff x=0$ \enspace et \enspace $e^{x}=2\iff x=\ln2$.\\
    $\res{S=\ens{0\pv\ln2}}$
  \item Avec $X=e^{x}>0$ : $(X-1)(X-2)\geq0$.\\
    Un trinôme est du signe de son coefficient dominant (ici $1>0$) à l'extérieur de ses racines, donc $X\leq1$ ou $X\geq2$.\\
    $e^{x}\leq1\iff x\leq0$ \enspace et \enspace $e^{x}\geq2\iff x\geq\ln2$.\\
    $\res{S=\intOF{-\infty}{0}\cup\intFO{\ln2}{+\infty}}$
\end{questions}
\end{minipage}\hfill
\begin{minipage}[t]{0.42\linewidth}\vspace{0pt}
\centering
\begin{tikzpicture}
\begin{axis}[otGraphe, width=\linewidth,
  xmin=-2.6, xmax=1.45, ymin=-0.75, ymax=3.25,
  xtick={-2,-1,1}, ytick={1,2,3}]
  \addplot[otGreen,line width=2.6pt,opacity=0.55] coordinates {(-2.6,0) (0,0)};
  \addplot[otGreen,line width=2.6pt,opacity=0.55] coordinates {(0.6931,0) (1.45,0)};
  \addplot[otMuted,dashed,line width=0.5pt,domain=-2.6:1.45] {2};
  \addplot[otNavy,line width=1.2pt,domain=-2.6:1.21,samples=200,smooth]
    {exp(2*x)-3*exp(x)+2};
  \addplot[only marks,mark=*,mark size=1.5pt,otNavy]
    coordinates {(0,0) (0.6931,0) (0.4055,-0.25)};
  \node[otEtiquette,anchor=north east] at (axis cs:-0.04,-0.09) {$0$};
  \node[otEtiquette,anchor=north west] at (axis cs:0.74,-0.09) {$\ln2$};
  \node[otEtiquette,anchor=south west,text=otMuted] at (axis cs:-2.55,2.04) {$y=2$};
  \node[anchor=east,font=\small,text=otNavy] at (axis cs:1.02,2.75) {$\Cf$};
\end{axis}
\end{tikzpicture}\\[1mm]
\legende{En vert : les solutions de $f(x)\geq0$.\\En pointillés : l'asymptote d'équation $y=2$.}
\end{minipage}

\begin{questions}[start=6]
  \item \methode{factoriser par le terme dominant}\\
    En $+\infty$ : pour tout réel $x$, $f(x)=e^{2x}\left(1-\dfrac{3}{e^{x}}+\dfrac{2}{e^{2x}}\right)$,
    et la parenthèse tend vers $1$. Comme $e^{2x}$ tend vers $+\infty$ :
    \[ \lim_{x\to+\infty}f(x)=\res{+\infty} \]
    En $-\infty$ : $e^{2x}$ et $e^{x}$ tendent vers $0$, donc
    \[ \lim_{x\to-\infty}f(x)=0-0+2=\res{2} \]
    La droite d'équation $y=2$ est asymptote horizontale à $\Cf$ en $-\infty$.
\end{questions}
\end{corrige}

%% ---------------------------------------------------------------------
\begin{corrige}{11}{Étude d'une fonction composée}
\noindent
\begin{minipage}[t]{0.54\linewidth}\vspace{0pt}\raggedright
\begin{questions}
  \item Avec $u(x)=x^{2}-2x$ : $h'(x)=(2x-2)e^{x^{2}-2x}=\res{2(x-1)e^{x^{2}-2x}}$, du signe de $x-1$. $h$ est décroissante sur
    $\intOF{-\infty}{1}$ et croissante sur $\intFO{1}{+\infty}$.
  \item $\res{\min h=h(1)=e^{-1}=\frac1e}$.
  \item $x^{2}-2x=x(x-2)\to+\infty$ en $\pm\infty$, et $e^{X}\to+\infty$ quand $X\to+\infty$ : $\res{\lim_{\pm\infty}h=+\infty}$.
  \item $(2-x)^{2}-2(2-x)=4-4x+x^{2}-4+2x=x^{2}-2x$, donc $h(2-x)=h(x)$ : la courbe est $\res{\text{symétrique par rapport à la droite }x=1}$.
\end{questions}
\end{minipage}\hfill
\begin{minipage}[t]{0.42\linewidth}\vspace{1mm}
\centering
\begin{tableau}
  \otTab[1.45]{$x$/0.8,$h'(x)$/0.8,$h$/1.6}{$-\infty$,$1$,$+\infty$}
  \tkzTabLine{,-,z,+,}
  \tkzTabVar{+/$+\infty$,-/$\frac1e$,+/$+\infty$}
\end{tableau}
\end{minipage}
\end{corrige}

%% ---------------------------------------------------------------------
\begin{corrige}{12}{Optimisation}
\noindent
\begin{minipage}[t]{0.54\linewidth}\vspace{0pt}\raggedright
\begin{questions}
  \item $k'(x)=2xe^{-x}-x^{2}e^{-x}=\res{x(2-x)e^{-x}}$ : positive sur $\intFF{0}{2}$, négative sur $\intFO{2}{+\infty}$.
  \item $\res{\max k=k(2)=\frac{4}{e^{2}}\approx0{,}541}$.
  \item $x^{2}e^{-x}=\frac{x^{2}}{e^{x}}\to\res{0}$ en $+\infty$ (croissances comparées), et $k(0)=0$.
\end{questions}
\end{minipage}\hfill
\begin{minipage}[t]{0.42\linewidth}\vspace{1mm}
\centering
\begin{tableau}
  \otTab[1.45]{$x$/0.8,$k'(x)$/0.8,$k$/1.6}{$0$,$2$,$+\infty$}
  \tkzTabLine{z,+,z,-,}
  \tkzTabVar{-/$0$,+/$\frac{4}{e^{2}}$,-/$0$}
\end{tableau}
\end{minipage}
\end{corrige}

%% ---------------------------------------------------------------------
\begin{corrige}{13}{Étude de convexité}
\noindent
\begin{minipage}[t]{0.54\linewidth}\vspace{0pt}\raggedright
\begin{questions}
  \item $\res{f'(x)=-2xe^{-x^{2}}}$, puis $f''(x)=-2e^{-x^{2}}+(-2x)(-2x)e^{-x^{2}}=\res{2\left(2x^{2}-1\right)e^{-x^{2}}}$.
  \item $f''(x)$ est du signe de $2x^{2}-1$ : $f$ est $\res{\text{convexe sur }\intOF{-\infty}{-\frac{\sqrt2}{2}}\text{ et sur }\intFO{\frac{\sqrt2}{2}}{+\infty}}$
    et $\res{\text{concave sur }\intFF{-\frac{\sqrt2}{2}}{\frac{\sqrt2}{2}}}$.
  \item $f''$ s'annule en changeant de signe : points d'inflexion $\res{\left(\pm\frac{\sqrt2}{2}\pv e^{-1/2}\right)}$, avec $e^{-1/2}\approx0{,}61$.
  \item $-x^{2}\to-\infty$, donc $\res{\lim_{\pm\infty}f=0}$ : l'axe des abscisses est asymptote. Le maximum est $f(0)=1$.
\end{questions}
\end{minipage}\hfill
\begin{minipage}[t]{0.42\linewidth}\vspace{1mm}
\centering
\begin{tikzpicture}
\begin{axis}[otSchema, axis lines=middle, width=68mm, height=40mm, xmin=-3.2, xmax=3.4, ymin=-0.1, ymax=1.2,
  xtick={-2,-1,1,2}, ytick={1}, xlabel={$x$}, xlabel style={at={(1,0)},anchor=west}]
  \addplot[otNavy,line width=1.2pt,domain=-3.1:3.1,samples=120] {exp(-x^2)};
  \addplot[only marks,mark=*,mark size=1.5pt,otGreenDark] coordinates {(-0.7071,0.6065) (0.7071,0.6065)};
\end{axis}
\end{tikzpicture}\\[0.5mm]
\legende{La courbe en cloche et ses deux points d'inflexion.}
\end{minipage}
\end{corrige}

%% ---------------------------------------------------------------------
\begin{corrige}{14}{La fonction sigmoïde}
\begin{questions}
  \item En $-\infty$, $e^{-x}\to+\infty$ donc $\res{\lim f=0}$ ; en $+\infty$, $e^{-x}\to0$ donc $\res{\lim f=1}$. Les droites $y=0$ et
    $y=1$ sont asymptotes à $\Cf$, respectivement en $-\infty$ et en $+\infty$.
  \item $f=\frac1v$ avec $v(x)=1+e^{-x}$ et $v'(x)=-e^{-x}$, donc $f'(x)=-\frac{v'(x)}{v(x)^{2}}=\res{\dfrac{e^{-x}}{\left(1+e^{-x}\right)^{2}}}$. Et
    $f(x)\big(1-f(x)\big)=\frac{1}{1+e^{-x}}\times\frac{e^{-x}}{1+e^{-x}}=f'(x)$. Comme $f'>0$, $\res{f\text{ est strictement croissante sur }\R}$.
\end{questions}
\noindent
\begin{minipage}[t]{0.52\linewidth}\vspace{0pt}\raggedright
\begin{questions}[start=3]
  \item $f(-x)=\frac{1}{1+e^{x}}=\frac{e^{-x}}{e^{-x}+1}$, d'où $f(-x)+f(x)=\frac{e^{-x}+1}{1+e^{-x}}=1$. Le milieu des points de $\Cf$
    d'abscisses $x$ et $-x$ est donc $\left(0\pv\frac12\right)$ : $\res{\Omega\text{ est centre de symétrie}}$.
  \item $f(0)=\frac12$ et $f'(0)=\frac12\times\frac12=\frac14$ : $\res{T:y=\frac x4+\frac12}$.
  \item $f'>0$ et $1-2f(x)>0\iff f(x)<\frac12\iff x<0$. $f$ est $\res{\text{convexe sur }\intOF{-\infty}{0}}$ et
    $\res{\text{concave sur }\intFO{0}{+\infty}}$ : $\Omega$ est un point d'inflexion. $\Cf$ est au-dessus de $T$ sur $\intOF{-\infty}{0}$ et
    en dessous sur $\intFO{0}{+\infty}$.
  \item $f(x)\geq0{,}9\iff1+e^{-x}\leq\frac{10}{9}\iff e^{-x}\leq\frac19\iff x\geq\ln9$ : $\res{S=\intFO{\ln9}{+\infty}}$, avec $\ln9\approx2{,}20$.
\end{questions}
\end{minipage}\hfill
\begin{minipage}[t]{0.44\linewidth}\vspace{1mm}
\centering
\begin{tikzpicture}
\begin{axis}[otSchema, axis lines=middle, width=72mm, height=46mm, xmin=-6.3, xmax=6.6, ymin=-0.12, ymax=1.25,
  xtick={-4,-2,2,4}, ytick={0.5,1}, yticklabels={$\frac12$,$1$}, xlabel={$x$}, xlabel style={at={(1,0)},anchor=west}]
  \draw[otGreenDark,dashed,line width=0.6pt] (axis cs:-6.2,1) -- (axis cs:6.2,1);
  \addplot[otGreen,line width=0.9pt,domain=-2.4:2.4] {x/4+0.5};
  \addplot[otNavy,line width=1.2pt,domain=-6.2:6.2,samples=120] {1/(1+exp(-x))};
  \fill[otNavy] (axis cs:0,0.5) circle (1.5pt);
  \node[anchor=north west,font=\scriptsize] at (axis cs:0.05,0.48) {$\Omega$};
  \node[anchor=west,font=\scriptsize,text=otGreenDark] at (axis cs:2.4,1.1) {$T$};
\end{axis}
\end{tikzpicture}\\[0.5mm]
\legende{La sigmoïde, sa tangente en $\Omega$ et ses asymptotes.}
\end{minipage}
\end{corrige}

%% ---------------------------------------------------------------------
\begin{corrige}{15}{Encadrer le nombre $e$}
\partieExo{A}{Une inégalité fondamentale}
\begin{questions}
  \item $\varphi(x)=e^{x}-1-x$ a pour dérivée $e^{x}-1$, négative sur $\intOF{-\infty}{0}$ et positive ensuite : son minimum est
    $\varphi(0)=0$, donc $\res{e^{x}\geq1+x}$ (la courbe de $\exp$ est au-dessus de sa tangente en $0$).
  \item $e^{1/n}\geq1+\frac1n>0$. La fonction $t\mapsto t^{n}$ étant croissante sur $\intFO{0}{+\infty}$ : $\left(e^{1/n}\right)^{n}=e\geq\left(1+\frac1n\right)^{n}$.
  \item $e^{-1/(n+1)}\geq1-\frac{1}{n+1}=\frac{n}{n+1}>0$. En passant aux inverses (la fonction inverse est décroissante sur
    $\intOO{0}{+\infty}$) : $e^{1/(n+1)}\leq\frac{n+1}{n}=1+\frac1n$, puis, à la puissance $n+1$ : $\res{e\leq\left(1+\frac1n\right)^{n+1}}$.
\end{questions}
\partieExo{B}{Deux suites qui encadrent $e$}
\begin{questions}[resume]
  \item $u_1=2$, $u_2=\frac94=2{,}25$ et $v_1=4$. Puis $v_n-u_n=\left(1+\frac1n\right)^{n}\left(1+\frac1n-1\right)=\res{\dfrac{u_n}{n}}$, et comme
    $0\leq u_n\leq e$, $\res{0\leq v_n-u_n\leq\frac en}$.
  \item $0\leq e-u_n\leq v_n-u_n\leq\frac en$, et $\frac en\to0$ : par encadrement, $\res{\lim u_n=e}$. Puis $v_n=u_n+\frac{u_n}{n}\to e$.
  \item Il suffit que $\frac en<10^{-2}$ : comme $e<2{,}72$, $\res{n=272}$ convient. On obtient $2{,}713\,3\leq e\leq2{,}723\,3$ : la
    convergence est lente, puisque l'écart ne décroît que comme $\frac1n$.
\end{questions}
\end{corrige}

%% ---------------------------------------------------------------------
\begin{corrige}{16}{Combien de solutions pour $e^{x}=kx$ ?}
\begin{questions}
  \item En $-\infty$ : $e^{x}\to0$ et $x\to-\infty$, donc $\res{\lim g=0}$. En $0^{-}$ : $e^{x}\to1$ et $x\to0^{-}$, donc $\res{\lim g=-\infty}$ ;
    en $0^{+}$, $\res{\lim g=+\infty}$. En $+\infty$ : $\res{\lim g=+\infty}$ par croissances comparées.
\end{questions}
\noindent
\begin{minipage}[t]{0.5\linewidth}\vspace{0pt}\raggedright
\begin{questions}[start=2]
  \item $g'(x)=\frac{e^{x}\times x-e^{x}\times1}{x^{2}}=\res{\dfrac{(x-1)e^{x}}{x^{2}}}$, du signe de $x-1$ : voir le tableau.
  \item $e^{0}=1\neq0=k\times0$. Pour $x\neq0$, on divise par $x$ : $e^{x}=kx\iff\frac{e^{x}}{x}=k$.
\end{questions}
\end{minipage}\hfill
\begin{minipage}[t]{0.46\linewidth}\vspace{1mm}
\centering
\begin{tableau}
  \otTab[1.2]{$x$/0.8,$g'(x)$/0.8,$g$/1.6}{$-\infty$,$0$,$1$,$+\infty$}
  \tkzTabLine{,-,d,-,z,+,}
  \tkzTabVar{+/$0$,-D+/$-\infty$/$+\infty$,-/$e$,+/$+\infty$}
\end{tableau}
\end{minipage}
\noindent
\begin{minipage}[t]{0.52\linewidth}\vspace{0pt}\raggedright
\begin{questions}[start=4]
  \item Sur $\intOO{-\infty}{0}$, $g$ est continue, strictement décroissante, de $0$ à $-\infty$ ; sur $\intOO{0}{+\infty}$, $g\geq e$. D'où :
    \begin{itemize}[label={\color{otNavy}\textbullet},leftmargin=1.2em,itemsep=0.3mm,topsep=0.5mm]
      \item si $k<0$ : $\res{\text{une solution}}$, négative ;
      \item si $0\leq k<e$ : $\res{\text{aucune solution}}$ ;
      \item si $k=e$ : $\res{\text{une solution}}$, $x=1$ ;
      \item si $k>e$ : $\res{\text{deux solutions}}$, l'une dans $\intOO{0}{1}$, l'autre dans $\intOO{1}{+\infty}$.
    \end{itemize}
  \item La tangente à la courbe de $\exp$ en $1$ a pour équation $y=e(x-1)+e=ex$ : elle passe par l'origine. Les solutions de
    $(E_k)$ sont les abscisses des points d'intersection de la courbe avec la droite $y=kx$ : au-delà de la pente $e$, la droite
    coupe la courbe deux fois ; en deçà (pente positive), elle passe sous la courbe convexe.
  \item Par balayage : $\res{0{,}61<\alpha<0{,}62}$ et $\res{1{,}51<\beta<1{,}52}$.
\end{questions}
\end{minipage}\hfill
\begin{minipage}[t]{0.44\linewidth}\vspace{1mm}
\centering
\begin{tikzpicture}
\begin{axis}[otGraphe, unit vector ratio=1 0.5, width=62mm, xmin=-1.8, xmax=2.2, ymin=-1.4, ymax=6.2, clip=true, clip mode=individual,
  xtick={-1,1,2}, ytick={1,2,3,4,5}]
  \addplot[otMuted,line width=0.7pt,domain=-1.8:2.2] {2*x};
  \addplot[otGreen,line width=0.9pt,domain=-1.8:2.2] {2.71828*x};
  \addplot[otOrange,line width=0.8pt,domain=-1.8:2.2] {3*x};
  \addplot[otNavy!60,line width=0.7pt,domain=-1.8:2.2] {-x};
  \addplot[otNavy,line width=1.3pt,domain=-1.8:1.85,samples=100] {exp(x)};
  \fill[otGreenDark] (axis cs:1,2.71828) circle (1.5pt);
  \node[anchor=south,font=\scriptsize,text=otOrange] at (axis cs:2.07,6.2) {$k=3$};
  \node[anchor=west,font=\scriptsize,text=otGreenDark] at (axis cs:2.22,5.98) {$k=e$};
  \node[anchor=west,font=\scriptsize,text=otMuted] at (axis cs:2.22,4.4) {$k=2$};
  \node[anchor=south west,font=\scriptsize,text=otNavy!70] at (axis cs:-1.75,2.0) {$k=-1$};
\end{axis}
\end{tikzpicture}\\[0.5mm]
\legende{La courbe de $\exp$ et quelques droites $y=kx$.}
\end{minipage}
\end{corrige}
